\par\indent\SourceLeaf{214}%
A triangulation network always extends over only a very limited part of the Earth’s surface. We can therefore replace \(f\upsilon=\upsilon\) with \(f\upsilon=\upsilon+\text{Const.}\) to extend the general solution a little further and achieve the stated purpose more fully. Giving this constant a real value would evidently gain us nothing. It would simply make \(T\) and \(t\) differ by this constant, changing only the origins from which longitude is measured. The situation is quite different if the constant is given an imaginary value. If we set it \(=-i\log k\), we obtain
\[
T=t,\qquad \mathrm{tang}\tfrac{1}{2}U=k\,\mathrm{tang}\tfrac{1}{2}w\mathord{.}(\frac{1+\varepsilon\cos w}{1-\varepsilon\cos w})^{\frac{1}{2}\varepsilon}
\]
\par\indent To decide on the most suitable value of \(k\), we must first determine the ratio of magnification.

In the notation of Articles 5 and 7, we have here
\[
\begin{aligned}
n&=aa\sin u^2\\
N&=AA\sin U^2\\
\varphi\upsilon&=1
\end{aligned}
\]
\par\indent Hence
\[
\begin{aligned}
m&=\frac{A\sin U}{a\sin u}=\frac{A\sin U}{a\sin w}\mathord{.}\BareRoot{(1-\varepsilon\varepsilon\cos w^2)}\\
&=\frac{A}{a}\mathord{.}\frac{k(1-\varepsilon\varepsilon\cos w^2)^{\frac{1}{2}+\frac{1}{2}\varepsilon}}{\cos\frac{1}{2}w^2(1-\varepsilon\cos w)^{\varepsilon}+kk\sin\frac{1}{2}w^2(1+\varepsilon\cos w)^{\varepsilon}}
\end{aligned}
\]
\noindent so this ratio depends only on latitude. We obtain the least possible departure from perfect similarity by choosing \(k\) so that \(m\) has equal values at the extreme latitudes. It follows that \(m\) at the middle latitude will be very close to its maximum or minimum. If we denote the extreme values of \(w\) by \(w^0\) and \(w'\), this gives
\[
k=\BareRoot{\frac{\displaystyle\frac{\cos\frac{1}{2}{w^0}^{2}(1-\varepsilon\cos w^0)^{\varepsilon}}{(1-\varepsilon\varepsilon\cos {w^0}^{2})^{\frac{1}{2}+\frac{1}{2}\varepsilon}}-\frac{\cos\frac{1}{2}{w'}^{2}(1-\varepsilon\cos w')^{\varepsilon}}{(1-\varepsilon\varepsilon\cos {w'}^{2})^{\frac{1}{2}+\frac{1}{2}\varepsilon}}}{\displaystyle\frac{\sin\frac{1}{2}{w'}^{2}(1+\varepsilon\cos w')^{\varepsilon}}{(1-\varepsilon\varepsilon\cos {w'}^{2})^{\frac{1}{2}+\frac{1}{2}\varepsilon}}-\frac{\sin\frac{1}{2}{w^0}^{2}(1+\varepsilon\cos w^0)^{\varepsilon}}{(1-\varepsilon\varepsilon\cos {w^0}^{2})^{\frac{1}{2}+\frac{1}{2}\varepsilon}}}}
\]
\noindent To find the latitude at which \(m\) reaches its greatest or least value, we have
\SourceLeaf{215}%
\[
\begin{aligned}
\frac{\mathrm{d}m}{m}&=\mathrm{cotang}\,U\mathord{.}\mathrm{d}u-\mathrm{cotang}\,w\mathord{.}\mathrm{d}w+\frac{\varepsilon\varepsilon\cos w\mathord{.}\sin w\mathord{.}\mathrm{d}w}{1-\varepsilon\varepsilon\cos w^2}\\
\frac{\mathrm{d}U}{\sin U}&=\frac{\mathrm{d}w}{\sin w}-\frac{\varepsilon\varepsilon\sin w\mathord{.}\mathrm{d}w}{1-\varepsilon\varepsilon\cos w^2}=\frac{(1-\varepsilon\varepsilon)\mathrm{d}w}{(1-\varepsilon\varepsilon\cos w^2)\sin w}
\end{aligned}
\]
\noindent and hence
\[
\frac{\mathrm{d}m}{m}=\frac{(1-\varepsilon\varepsilon)\mathrm{d}w}{\sin w(1-\varepsilon\varepsilon\cos w^2)}\mathord{.}(\cos U-\cos w)
\]
\noindent This shows that \(m\) reaches its greatest or least value when \(U=w\). If we denote the value of \(w\) at that point by \(W\), we get
\[
k=(\frac{1-\varepsilon\cos W}{1+\varepsilon\cos W})^{\frac{1}{2}\varepsilon}\quad\text{or}\quad \cos W=\frac{1-k^{\frac{2}{\varepsilon}}}{\varepsilon(1+k^{\frac{2}{\varepsilon}})}
\]
\noindent from which we can determine \(W\) once \(k\) has been calculated by the formula above. In practice, exact equality of the values of \(m\) at the extreme latitudes is of little importance. We can simply choose approximately the middle latitude for \(90^0-W\) and derive \(k\) from it. The general relation between \(U\) and \(w\) is then given by the formula
\[
\mathrm{tang}\tfrac{1}{2}U=\mathrm{tang}\tfrac{1}{2}w\left\{\frac{(1-\varepsilon\cos W)(1+\varepsilon\cos w)}{(1+\varepsilon\cos W)(1-\varepsilon\cos w)}\right\}^{\frac{1}{2}\varepsilon}
\]
\noindent For actual numerical calculation, however, it is more advantageous to use series. They can be put in various forms, and we will leave their derivation aside here.

It is easy to see that when \(w<W\), we have \(U>w\), so \(\cos U-\cos w\), and consequently also \(\frac{\mathrm{d}m}{\mathrm{d}w}\), is negative. When \(w>W\), we have \(U<w\), and hence \(\frac{\mathrm{d}m}{\mathrm{d}w}\) is positive. It is therefore clear that when \(w=U=W\), the value of \(m\) is always a minimum, namely
\[
=\frac{A}{a}\BareRoot{(1-\varepsilon\varepsilon\cos W^2)}
\]
\noindent Thus, if we choose the radius of the sphere \(A=\frac{a}{\BareRoot{(1-\varepsilon\varepsilon\cos W^2)}}\), the representation of infinitely small parts of the ellipsoid at latitude \(90^0-W\) is similar and indeed equal to the original. At other latitudes it is larger.

It is advantageous to expand the logarithm of \(m\) in a series in successive powers of \(\cos U-\cos W\). Its first terms, sufficient for practical use, are
\SourceLeaf{216}%
\[
\begin{aligned}
\log\mathrm{hyp}.\,m&=\log\{\frac{A}{a}\BareRoot{(1-\varepsilon\varepsilon\cos W^2)}\}+\frac{\varepsilon\varepsilon}{2(1-\varepsilon\varepsilon)}\mathord{.}(\cos U-\cos W)^2\\
&\qquad-\frac{2\varepsilon^4\cos W}{3(1-\varepsilon\varepsilon)^2}\mathord{.}(\cos U-\cos W)^3\ .\ .\ .\ .
\end{aligned}
\]
\par\indent If, for example, we transfer the Danish monarchy between the latitude limits \(53^0\) and \(58^0\) to the sphere in this way and set \(W=34^0\,30'\), then with flattening \(\frac{1}{303}\) the representation at the limits is enlarged by only \(\frac{1}{530000}\) in its linear dimensions.

We must content ourselves here with giving only a brief indication of \textit{one} use of transferring figures in higher geodesy, and reserve a proper treatment for another place.

